\documentclass{article}

\usepackage{hyperref}
\usepackage{amsmath}
\usepackage{amssymb}
\usepackage{amsthm}
\usepackage{dsfont}
\usepackage{stmaryrd}
\usepackage{xcolor}
\usepackage{tikz}
\usepackage{cleveref}
\usepackage{fullpage}
\usepackage{authblk}

% Helps reduce minor overfull \hbox warnings.
\usepackage{microtype}

\title{Formalizing Flag Algebra in Lean}
\date{}

\author[1]{Gyeongwon Jeong}
\affil[1]{School of Computing, KAIST, South Korea,
  \texttt{jgyw0910@kaist.ac.kr}}

\author[2]{Seonghun Park}
\affil[2]{School of Computing, KAIST, South Korea,
  \texttt{hun57@kaist.ac.kr}}

\author[3]{Hongseok Yang}
\affil[3]{School of Computing, KAIST, South Korea,
  \texttt{hongseok00@gmail.com}}

\begin{document}
\maketitle

\begin{abstract}
TODO: Write an abstract.
\end{abstract}

\section{Introduction}

\section{Recommended Papers}

Lean formalization papers recommended by Gemini from \url{https://leanprover-community.github.io/papers.html}.

\subsection{\href{https://drops.dagstuhl.de/storage/00lipics/lipics-vol237-itp2022/LIPIcs.ITP.2022.9/LIPIcs.ITP.2022.9.pdf}{\textbf{Formalising Szemerédi's Regularity Lemma in Lean}} (Yaël Dillies and Bhavik Mehta, 2022)}

\subsubsection{Contributions of the Paper}
\begin{itemize}
    \item \textbf{Equitable Regularity Lemma:} Formalized a stronger, previously unattempted version of the lemma requiring parts of approximately equal size.
    \item \textbf{Explicit Upper Bound:} Derived a rigorous, explicit upper bound for the number of partitions directly within the Lean code.
    \item \textbf{Theorem Applications:} Applied the lemma to successfully formalize the Triangle Counting Lemma, Triangle Removal Lemma, Roth's theorem, and the Corners theorem.
\end{itemize}

\subsubsection{Structure of the Paper}
\begin{itemize}
    \item \textbf{Section 1 \& 2 (Background \& Math):} Introduces the mathematical context and provides pen-and-paper proof sketches before introducing code.
    \item \textbf{Section 3 \& 4 (Core Formalization):} Details the Lean implementation of the Regularity Lemma and Triangle lemmas, covering definitions, witness handling, and bounds calculations.
    \item \textbf{Section 5 (Extensions):} Explains the construction of tripartite graphs and definitional workarounds needed for the Corners theorem and Roth's theorem.
    \item \textbf{Section 6 (Discussion):} Compares the methodology with a concurrent Isabelle/HOL study and outlines future applications.
\end{itemize}

\subsubsection{Key Takeaways for Writing Formalization Papers}
\begin{itemize}
    \item \textbf{Structure by Abstraction:} Presenting mathematical proof sketches before technical code significantly reduces reader fatigue.
    \item \textbf{Defend Design Choices:} Clearly articulate the rationale behind creating custom definitions (e.g., \texttt{finpartition}) instead of relying solely on existing library structures.
    \item \textbf{Transparent Limitations:} Openly acknowledge the theorem prover's constraints (e.g., limitations with natural number subtraction or non-linear arithmetic) and explain the specific workarounds used.
    \item \textbf{Technical Benchmarking:} Compare related formalization work using precise technical metrics (e.g., differing definitions of uniformity) rather than general claims of novelty.
\end{itemize}

\subsection{\href{https://link.springer.com/chapter/10.1007/978-3-031-16681-5_5}{\textbf{Formalising the Kruskal-Katona Theorem in Lean}} (Bhavik Mehta, 2022)}

\subsubsection{Contributions of the Paper}
\begin{itemize}
    \item \textbf{Core Combinatorial Theorems:} Successfully formalized the Kruskal-Katona theorem, along with Sperner's theorem, the LYM inequalities, and the Erd\H{o}s-Ko-Rado theorem.
    \item \textbf{Computable Operators:} Implemented a computable version of the 'shadow' operator for finite sets, allowing for direct evaluation and automated testing within the prover.
    \item \textbf{Foundational Structures:} Formalized the colexicographic ordering on finite sets and the mathematical technique of UV-compressions.
\end{itemize}

\subsubsection{Structure of the Paper}
\begin{itemize}
    \item \textbf{Sections 1 \& 2 (Context):} Introduces the history of extremal combinatorics and discusses related formalization efforts in Isabelle and Lean.
    \item \textbf{Section 3 (Mathematical Background):} Provides informal, pen-and-paper proof sketches for the LYM inequalities, Kruskal-Katona, and Erd\H{o}s-Ko-Rado theorems to ground the reader.
    \item \textbf{Section 4 (Formalizing the Shadow):} Details the Lean definitions for set families and shadows, followed by the formal proofs of the local and general LYM inequalities.
    \item \textbf{Section 5 (Kruskal-Katona):} Explains the implementation of colexicographic ordering, the complex manual set calculations required for compressions, and the main theorem via strong induction.
    \item \textbf{Section 6 (Discussion):} Concludes by highlighting the wider applicability of compression techniques in extremal combinatorics.
\end{itemize}

\subsubsection{Key Takeaways for Writing Formalization Papers}
\begin{itemize}
    \item \textbf{Highlight Computability:} Designing definitions to be explicitly computable (e.g., using \texttt{finset.bind} and \texttt{image}) allows for direct sanity checks via \texttt{\#eval}. Showcasing this bridges the gap between abstract math and software engineering.
    \item \textbf{Strategic Generalization:} Start by defining concepts over general, arbitrary finite types, and specialize to specific sets (like \texttt{fin n}) only when absolute metrics (like family measure) are required.
    \item \textbf{Document Automation Gaps:} Transparently discuss where theorem prover automation falls short (e.g., handling lengthy set equality calculations involving singletons) and explain the manual chains of elementary lemmas required as workarounds.
    \item \textbf{Adapt Mathematical Proofs:} Demonstrate how traditional textbook proofs (like double-counting on pairs) are translated into structures better suited for the prover, such as forming two exact sets and demonstrating subset relations.
\end{itemize}

\subsection{\href{https://drops.dagstuhl.de/storage/00lipics/lipics-vol309-itp2024/LIPIcs.ITP.2024.35/LIPIcs.ITP.2024.35.pdf}{\textbf{Formal Verification of the Empty Hexagon Number}} (Bernardo Subercaseaux et al., 2024)}

\subsubsection{Contributions of the Paper}
\begin{itemize}
    \item \textbf{Formal Verification of the SAT Reduction:} Rigorously verified in Lean the process of converting the theorem (that any set of 30 points contains an empty convex hexagon, $h(6) \le 30$) into a SAT problem.
    \item \textbf{Equivalence between Geometry and Combinatorics:} Mathematically proved that continuous geometric definitions (e.g., convex hulls, points in triangles) are equivalent to discrete combinatorial abstractions (Triple Orientations).
    \item \textbf{Verification of Symmetry Breaking:} Formally proved that transforming the point sets into a 'canonical position' to reduce the search space does not result in a loss of generality.
\end{itemize}

\subsubsection{Structure of the Paper}
\begin{itemize}
    \item \textbf{Sections 1 \& 2 (Background \& Overview):} Introduces the history of the Erd\H{o}s-Szekeres problem and the overall pipeline of the SAT-based proof.
    \item \textbf{Sections 3 \& 4 (Geometry \& Discretization):} Covers the process of discretizing the continuous problem space using geometric definitions from \texttt{mathlib} and determinant-based triple orientations.
    \item \textbf{Sections 5 \& 6 (Symmetry Breaking \& SAT Encoding):} Explains the search space reduction techniques via canonical position transformations, the construction of the CNF formula, and the actual execution results of the SAT solver and proof checker (\texttt{cake\_lpr}).
    \item \textbf{Sections 7 \& 8 (Related Work \& Conclusion):} Concludes by providing a concrete comparison of encoding sizes and performance metrics against prior Isabelle-based $g(6)$ verification studies.
\end{itemize}

\subsubsection{Key Takeaways for Writing Formalization Papers}
\begin{itemize}
    \item \textbf{Connect to Established Libraries:} Rather than axiomatizing new concepts from scratch, link them to the foundational definitions of standard libraries (such as \texttt{mathlib}'s convex hulls and determinants) to increase the overall trustworthiness of the system.
    \item \textbf{Verify the Reduction Gap:} When combinatorially encoding a complex mathematical problem so that external computing tools (like SAT solvers) can solve it, clearly document how the validity of the translation logic itself was secured within the theorem prover.
    \item \textbf{Clear Methodological Comparison:} When comparing against prior work, do not simply state that a different tool was used. Prove superiority by providing explicit metrics, such as the encoding formula size ($O(n^4)$ vs. $O(n^6)$) and actual verification time (e.g., CPU hours).
    \item \textbf{Modular Trust Story:} For proofs requiring massive computation, it is crucial to clearly explain the "trust structure"—detailing how the theorem prover (Lean), high-performance solver (CaDiCaL), and external checker (\texttt{cake\_lpr}) are separated and integrated so that the entire pipeline operates without logical flaws.
\end{itemize}

\subsection{\href{https://arxiv.org/abs/2101.00127}{\textbf{Formalizing Hall's Marriage Theorem in Lean}} (Alena Gusakov, Bhavik Mehta, and Kyle Miller, 2021)}

\subsection{\href{https://drops.dagstuhl.de/storage/00lipics/lipics-vol141-itp2019/LIPIcs.ITP.2019.15/LIPIcs.ITP.2019.15.pdf}{\textbf{Formalizing the Solution to the Cap Set Problem}} (Sander Dahmen, Johannes Hölzl, and Robert Lewis, 2019)}


\end{document}
